\documentclass[10pt,a4paper]{amsart}
\usepackage[margin=19mm]{geometry}
\usepackage{amsmath,amssymb,mathtools}
\usepackage[scr=rsfs,cal=euler]{mathalfa}
\usepackage{xfrac}
\usepackage[unicode,hidelinks,bookmarks=false]{hyperref}
\newcommand{\ie}{i.e.\ }
\newcommand{\cf}{cf.\ }
\newcommand{\resp}{respectively\ }
\newcommand{\Subsection}{\subsection}
\providecommand{\nobreakdash}{\nobreak\hbox{-}}
\setlength{\emergencystretch}{2em}
\multlinegap0pt
\usepackage{multirow}
\usepackage{amssymb}
\usepackage{xcolor}
\usepackage{booktabs}
\usepackage{algorithm}\usepackage{algpseudocode}
\usepackage[all]{xy}
\newtheorem{proposition}{Proposition}
\title{Quantization as a Categorical Equivalence for Hilbert Bimodules and Lagrangian Relations}
\author{Benjamin H. Feintzeig}
\date{Source: arXiv:2602.15188v1 (CC BY 4.0)}
\begin{document}
\maketitle

\noindent\emph{Source and licence.} Quantization as a Categorical Equivalence for Hilbert Bimodules and Lagrangian Relations. Version arXiv:2602.15188v1 (CC BY 4.0), distributed by arXiv under the Creative Commons Attribution 4.0 International licence, http://creativecommons.org/licenses/by/4.0/, as recorded in the arXiv licence metadata for this exact version and shown on the abstract page https://arxiv.org/abs/2602.15188v1. Full licence notice: \texttt{licenses/arxiv-2602.15188-CC-BY-4.0.txt}.\par\smallskip

\noindent\textit{Editorial scope.} Counted unit: the article's proposition proposition (source0.tex lines 743--745). The article's complete original proof of it is delivered with it (source0.tex lines 747--765, one \texttt{proof} environment of 19 lines). No mathematics is written, repaired or re-derived: the statement and the proof are the article's own lines, in the article's own notation, and no retained line is substituted or re-flowed.  No further source-local context is needed: every cross-reference of every retained line resolves inside the two delivered spans.  Nothing was omitted from any retained span, so the package carries no omission record.  No retained line contains a citation, so the wrapper carries no bibliography and no bibliography entry.  No retained line contains a figure environment, an image-inclusion command or drawing code, so no image is delivered and the package declares no asset dependency.  Editorial formatting only, in addition: an AMS \texttt{amsart} preamble that restates the article's own declarations and macros used by the retained lines, namely the environments \texttt{proposition} on separate counters, together with the standard packages those lines need.  The retained environments print in this document as Proposition 1 in order of appearance, whereas the article prints the same items with the numbering of its own full document.  The complete native source of the article as distributed in the arXiv e-print for this exact version is bundled unchanged as \texttt{sources/194779/source0.tex}.
\begin{proposition}
    The map $\iota^*_q$ is an isomorphism $T^*_qM\oplus \mathcal{I}_q^*\cong\mathfrak{G}^*_{t_M(q)}\oplus\mathfrak{H}^*_{s_M(q)}$.
\end{proposition}

\begin{proof}
    Define the map $(\iota^*_q)^{-1}: \mathfrak{G}_{t_M(q)}^*\oplus \mathfrak{H}_{s_M(q)}^*\to T^*_qM\oplus\mathcal{I}_q^*$ by
    \begin{align}
        (\iota^*_q)^{-1}(\Phi)(\xi_q\oplus Z) = \Phi(\iota_q^{-1}(\xi_q\oplus Z))
    \end{align}
for $\Phi\in \mathfrak{G}_{t_M(q)}^*\oplus \mathfrak{H}_{s_M(q)}^*$, $\xi_q\in T_qM$, and $Z\in \mathcal{I}_q$.  Then $(\iota_q^*)^{-1}$ is indeed the inverse of $\iota_q^*$ because for all $\eta_q\in T_q^*M$, $\Phi\in \mathcal{I}_q^*$, $\xi_q\in T_qM$, and $Z\in \mathcal{I}_q$, we have
\begin{align}
    \big((\iota_q^*)^{-1}\circ\iota_q^*\big)(\eta_q\oplus\Phi)(\xi_q\oplus Z) &= \iota_q^*(\eta_q\oplus \Phi)(\iota_q^{-1}(\xi_q\oplus Z))\\
    &= (\eta_q\oplus \Phi)(\iota_q\circ\iota_q^{-1}(\xi_q\oplus Z))\\
    &=(\eta_q\oplus \Phi)(\xi_q\oplus Z),
\end{align}
and for all $\Phi\in \mathfrak{G}_{t_M(q)}^*\oplus \mathfrak{H}_{s_M(q)}^*$ and $Z\in \mathfrak{G}_{t_M(q)}\oplus\mathfrak{H}_{s_M(q)}$, we have
\begin{align}
    \big(\iota^*_q\circ(\iota_q^*)^{-1}\big)(\Phi)(Z) &= (\iota_q^*)^{-1}(\Phi)(\iota_q(Z))\\
    &= \Phi(\iota_q^{-1}\circ\iota_q(Z))\\
    &= \Phi(Z).
\end{align}

\end{proof}

\end{document}
